\documentclass[14pt]{article}
\usepackage[top=25truemm,bottom=20truemm,left=25truemm,right=25truemm]{geometry}
\usepackage{xskak}
\usepackage{fancyhdr}
\setlength\parindent{0pt}

\begin{document}




\newpage
\textbf{\Large Q. Show a winning sequence.}\\
{\setchessboard{boardfontsize=30pt,labelfontsize=14pt}
\newchessgame
\chessboard[setfen=q5r1/8/8/8/8/8/Q4K1k/1R6 w - - 0 1]}\\

\textbf{\Large Answer:}\\
\variation{\Large 1.Rh1+! Qxh1 2.Qxg3}\\
\variation{\Large 1.Rh1+! Kxh1 2.Qb1+ Rg1 3.Qh7#}\\

\vfill\textbf{\Large Source: }


\newpage
\noindent\textbf{\Large Q2. Is there any difference between left and right?}\\
{\setchessboard{boardfontsize=20pt,labelfontsize=12pt}
\newchessgame
\chessboard[setfen=6k1/7p/7K/6P1/8/8/7P/8 w - - 0 1]}
{\setchessboard{boardfontsize=20pt,labelfontsize=12pt}
\newchessgame
\chessboard[setfen=6k1/7p/7K/6P1/8/7P/8/8 w - - 0 1]}\\

\textbf{\Large Answer:}\\
\textbf{\Large White wins in left and draws in right.}\\
White wins as follow.\\
\variation{\Large 1.h3 Kh8}\\
\variation{\Large 2.h4 Kg8}\\
\variation{\Large 3.h5 Kh8}\\
\variation{\Large 4.g6 hxg6}\\
\variation{\Large 5.hxg6 Kg8}\\
\variation{\Large 6.g7}\\
\variation{\Large 1.h3 Kh8 2.h4 Kg8 3.h5 Kh8 4.g6 hxg6 5.hxg6 Kg8 6.g7}\\

\vfill\textbf{\Large Source: Theoretical endgame}




\newpage
\noindent\textbf{\Large Q. Mate in 3.}\\
{\setchessboard{boardfontsize=30pt,labelfontsize=14pt}
\newchessgame
\chessboard[setfen=7B/bB2PbR1/4P1p1/4rPN1/1P1kP3/4N3/3K4/4n3 w - - 0 1]}\\

\textbf{\Large Answer:}\\
\variation{\Large 1. Bd5!}\\
Any moves except e5 Rook, \variation{\Large 2. Nc2+ Nxc2 3. Nf3#}.\\

\vfill\Large {Source: Valentin Fedorovich Rudenko,  Problemnoter, 1960 (YACDB: 57240)}


\newpage
\noindent\textbf{\Large Q. Mate in 3.}\\
{\setchessboard{boardfontsize=30pt,labelfontsize=14pt}
\newchessgame
\chessboard[setfen=7B/bB2PbR1/4P1p1/4rPN1/1P1kP3/4N3/3K4/4n3 w - - 0 1]}\\

\textbf{\Large Answer:}\\
\variation{\Large 1. Bd5!}\\
Any moves except e5 Rook, \variation{\Large 2. Nc2+ Nxc2 3. Nf3#}.\\

\vfill\Large {Source: Valentin Fedorovich Rudenko,  Problemnoter, 1960 (YACDB: 57240)}



\newpage
\noindent\textbf{\Large Q. Mate in 2.}\\
{\setchessboard{boardfontsize=30pt,labelfontsize=14pt}
\newchessgame
\chessboard[setfen=8/4Q3/1qp1N2b/1p3n2/2B5/2BNk1p1/6P1/R3K2R w KQ - 0 1]}\\

\textbf{\Large Answer:}\\
\variation{\Large 1. O-O-O}\\

\vfill\Large {Source: Michael Lipton, Great Britain - Israel, 1961 (YACDB: 36663)}


32371
Mansfield, Comins
La Settimana Enigmistica, 1935
2nd Prize, Categoria I, 1st half-year
La Settimana Enigmistica, 1935

32806
Morse, Christopher Jeremy
The Observer, 1962


\newpage
\textbf{\Large Q1A. Mate in two. (4 marks)}\\
{\setchessboard{boardfontsize=30pt,labelfontsize=14pt}
\newchessgame
\chessboard[setfen=8/4K3/8/4k3/4B3/6R1/8/4Q3 w - - 0 1]}\\

\textbf{\Large Answer:}\\

\variation{\Large 1. Re3! Kf4 2. Qg3\#}\\
\variation{\Large 1... Kd4 2. Bc3\#}\\

\vfill\Large{Source: Smuel Loyd, Household Journal, 1861}
%% Intention: Symmetry, zugzang (Easy warm-up)


\newpage
\textbf{\Large Q1B. Mate in two. (4 marks)}\\
{\setchessboard{boardfontsize=30pt,labelfontsize=14pt}
\newchessgame
\chessboard[setfen=2B5/8/8/2K3P1/4k3/6R1/3B4/8 w - - 0 1]}\\

\textbf{\Large Answer:}\\

\variation{\Large 1. Rg2! Kf3 2. Bb7\#}\\
\variation{\Large 1... Kd3 2. Bf5\#}\\
\variation{\Large 1... Ke5 2. Re2\#}\\

\vfill\Large{Source: Smuel Loyd, Musical World, 1859}
%% Intention: Two bishops (Easy warm-up) The previous one or this one for easy warm-up.



\newpage
\textbf{\Large Q3. Mate in 2.}\\
{\setchessboard{boardfontsize=30pt,labelfontsize=14pt}
\newchessgame
\chessboard[setfen=K1Q5/1p5r/8/2B5/8/8/R7/7k w - - 0 1]}\\

\textbf{\Large Answer:}\\
\variation{\Large 1. Bd4} creates another threats Qc1 in addition to Qh3 and Qh8.\\
\variation{\Large 1... Rd7 2. Qc1#}\\
\variation{\Large 1... Rc7 2. Qh3#}\\
Other bishop moves fails, as \variation{\Large 1. Bd4} is necessary to prevents Black Rook to d1 square.\\

\vfill{\Large Source: Miroslav Subotić, The Problemist, 1992 (YACPDB: 13286)}



\end{document}